L'apprentissage par renforcement (RL) est un domaine riche, mais dont la mise en pratique se heurte à un défi récurrent dès lors que l'environnement se complexifie : la rareté du signal de récompense (\textit{sparse reward}). Lorsque l'objectif final d'une tâche n'est atteint qu'occasionnellement -- voire jamais durant les premières phases de l'entraînement --, l'agent ne reçoit quasiment aucun signal exploitable pour orienter l'amélioration de sa politique, ce qui peut se traduire par une absence totale d'apprentissage observable, même après un nombre considérable de pas d'entraînement. Ce mémoire a pour objectif d'explorer concrètement ce défi, en progressant d'environnements standards à récompense dense vers un environnement avancé structurellement dépendant d'un signal rare.

Le premier axe concerne le problème de sélection et d'optimisation des modèles. Face à un problème donné, plusieurs décisions cruciales s'imposent : comment choisir d'abord le bon algorithme (DQN, PPO, SAC, etc.) en fonction des spécificités de l'environnement (comme des espaces d'actions discrets ou continus) ? Ensuite, comment relever le défi de l'optimisation des hyperparamètres, étape indispensable pour faire converger ces modèles au sein d'un environnement donné ?

Le second axe se concentre sur le problème du sparse reward lui-même : comment construire un signal de récompense suffisamment informatif pour guider l'apprentissage d'un agent lorsque l'objectif final demeure rare, sans pour autant dénaturer la tâche initiale ? Cette question soulève elle-même un risque supplémentaire, largement documenté dans la littérature : celui de comportements opportunistes inattendus (\textit{reward hacking}), lorsque le signal construit pour pallier la rareté de la récompense se révèle, en pratique, imparfaitement spécifié.

Ainsi, la problématique centrale de ce travail est de déterminer comment mener à bien l'apprentissage par renforcement dans un environnement à récompense rare : de la sélection et du réglage des algorithmes sur des cas d'étude standard à récompense dense, jusqu'à la conception d'un signal de récompense capable de surmonter la rareté structurelle de l'objectif final dans un environnement avancé comme Rocket League.